
\chapter{ORION - Ocena Ryzyka Informatycznego Organizacji Nadzorowanych}

\section{Luka badawcza}
\label{sec:luka}

Istniejące metodologie obejmują zarządzanie bezpieczeństwem, zarządzanie ryzykiem i~kontrole organizacyjne w znacznym stopniu w kierunku analizy systemów informatycznych. Są to standardy rynkowe, takie jak metodologie zdefiniowane w  ISO/IEC 27001, NIST Cybersecurity Framework, COBIT. Jednak żadna z nich nie spełnia jednocześnie  zasadniczych wymagań stawianych przez współczesne środowiska regulacyjne. Po pierwsze, brakuje im wielostronnej perspektywy obejmującej cały łańcuch dostaw ICT, od nadzorowanego podmiotu, poprzez jego bezpośrednich dostawców, po podprocesory działające na wielu poziomach. Po drugie, nie zawierają one opartego na dowodach mechanizmu punktacji, który rozróżniałby zweryfikowaną zgodność od zwykłego zapewnienia deklaratywnego rozróżnienia, które jest fundamentalne dla przeciwdziałania dobrze udokumentowanemu zjawisku \textit{teatru zgodności}. Po trzecie, nie zapewniają one wieloaspektowej struktury analitycznej, która formalnie przypisuje pytania audytowe do odpowiedniego podmiotu w organizacji, tj. radcy prawnego, biznesu, architekta oprogramowania lub administratora infrastruktury, w oparciu o charakter pytania, a nie wygodę organizacji. Po czwarte, nie oferują modułowej rozszerzalności sektorowej, która zachowuje uniwersalny rdzeń metodologiczny, jednocześnie aktywując perspektywy i~zestawy pytań specyficzne dla danej dziedziny dla regulowanych branż, takich jak finanse czy opieka zdrowotna.

\renewcommand{\arraystretch}{1.4}
% Use longtable instead of table + tabular
% Note: Total widths set to approx 16.5cm to fit A4
\begin{longtable}{p{4.2cm}p{2.2cm}p{2.2cm}p{2.2cm}p{2.4cm}p{1.8cm}}
\caption{Analiza porównawcza frameworków stosowanych w polskiej praktyce audytowej względem wymagań ORION} \label{tab:porownanie-frameworkow} \\
\toprule
\textbf{Wymaganie} 
& \textbf{ISO/IEC 27001} 
& \textbf{NIST CSF} 
& \textbf{COBIT 2019} 
& \textbf{Narzędzia GRC} 
& \textbf{ORION} \\
\midrule
\endfirsthead

% This header appears on subsequent pages
\multicolumn{6}{c}%
{{\bfseries \tablename\ \thetable{} -- kontynuacja z poprzedniej strony}} \\
\toprule
\textbf{Wymaganie} 
& \textbf{ISO/IEC 27001} 
& \textbf{NIST CSF} 
& \textbf{COBIT 2019} 
& \textbf{Narzędzia GRC} 
& \textbf{ORION} \\
\midrule
\endhead

\midrule
\multicolumn{6}{r}{{Kontynuacja na następnej stronie...}} \\
\endfoot

\bottomrule
\endlastfoot

\textbf{Pokrycie całego łańcucha dostaw ICT} 
& Częściowe\newline(tylko załącznik dot.\ dostawców) 
& Częściowe\newline(profil łańcucha dostaw) 
& Częściowe\newline(zarządzanie dostawcami) 
& Częściowe\newline(zależne od narzędzia) 
& \textbf{Pełne} \\
\midrule

\textbf{Ocena oparta na dowodach} 
& Nie 
& Nie 
& Nie 
& Częściowe\newline(zależne od konfiguracji) 
& \textbf{Tak} \\
\midrule

\textbf{Wieloperspektywiczne przypisanie pytań}\newline(prawna / biznesowa / techniczna) 
& Nie 
& Nie 
& Częściowe\newline(tylko macierz RACI) 
& Nie 
& \textbf{Tak} \\
\midrule

\textbf{Formalny model grafowy zależności ICT} 
& Nie 
& Nie 
& Nie 
& Nie 
& \textbf{Tak} \\
\midrule

\textbf{Propagacja ryzyka przez warstwy łańcucha dostaw} 
& Nie 
& Nie 
& Nie 
& Nie 
& \textbf{Tak} \\
\midrule

\textbf{Scoring zgodności z regulacjami}\newline(DORA / NIS2 / RODO) 
& Częściowe\newline(istnieją mapowania) 
& Częściowe\newline(istnieją mapowania) 
& Częściowe\newline(istnieją mapowania) 
& Częściowe\newline(zależne od konfiguracji) 
& \textbf{Tak} \\
\midrule

\textbf{Rozróżnienie wymogów obowiązkowych od dobrych praktyk} 
& Nie 
& Nie 
& Nie 
& Częściowe 
& \textbf{Tak} \\
\midrule

\textbf{Modułowa rozszerzalność sektorowa}\newline(finanse, zdrowie, energetyka) 
& Nie 
& Częściowe\newline(profile branżowe) 
& Nie 
& Częściowe\newline(płatne moduły) 
& \textbf{Tak} \\
\midrule

\textbf{Penalizacja niezweryfikowanych deklaracji} 
& Nie 
& Nie 
& Nie 
& Nie 
& \textbf{Tak} \\
\midrule

\textbf{Identyfikacja dostawców kluczowych} 
& Częściowe 
& Częściowe 
& Częściowe 
& Częściowe 
& \textbf{Tak} \\
\midrule

\textbf{Ocena planu wyjścia (exit plan)} 
& Nie 
& Nie 
& Częściowe 
& Częściowe 
& \textbf{Tak} \\
\midrule

\textbf{Transparentna, replikowalna metodyka} 
& Tak 
& Tak 
& Tak 
& Nie\newline(metodyka zastrzeżona) 
& \textbf{Tak} \\
\midrule

\textbf{Formalny model matematyczny} 
& Nie 
& Nie 
& Nie 
& Nie 
& \textbf{Tak} \\
\midrule

\textbf{Klasyfikacja ryzyka wg pochodzenia geograficznego}\newline(EOG / USA / pozostałe) 
& Nie 
& Nie 
& Nie 
& Nie 
& \textbf{Tak} \\

\end{longtable}

\begin{flushleft}
\footnotesize
\textit{Legenda:} \textbf{Tak} - wymaganie w pełni adresowane; 
\textbf{Częściowe} - wymaganie adresowane w ograniczonym zakresie lub wymagające dodatkowej konfiguracji; 
\textbf{Nie} - wymaganie nieadresowane.\\
\textit{Uwaga:} Porównanie obejmuje frameworki powszechnie stosowane w polskiej praktyce audytowej 
i compliance. Frameworki NIST SP~800-161 oraz ISACA CMMI, choć obecne w literaturze akademickiej, 
nie są szeroko stosowane jako samodzielne metodyki przez polskie organizacje 
i zostały pominięte w niniejszym zestawieniu.\\
\textit{Skróty:} CSF - \textit{Cybersecurity Framework}; 
GRC - \textit{Governance, Risk and Compliance} (ład, ryzyko i~zgodność).

\end{flushleft}



Mimo intensywności regulacyjnej prowadzonej przez regulatorów krajowych oraz międzynarodowych, brakuje usystematyzowanej metodyki, która pozwoliłaby organizacjom w sposób transparentny, co bardzo istotne powtarzalny i~oparty na obiektywnych dowodach ocenić  stan zgodności wdrożonych systemów informatycznych wraz z ich łańcuchem dostawców. Dobrze przeprowadzony audyt systemów tworzy raport czytelny i~zrozumiały dla wszystkich interesariuszy i~faktycznie wskazuje na niezgodności.
Istniejące narzędzia i~frameworki, norma ISO/IEC 27001,  NIST Cybersecurity Framework, model COBIT czy metodyki ISACA, dostarczają  wskazówek w zakresie zarządzania bezpieczeństwem informacji, jednak żaden z nich nie adresuje jednocześnie: 
\begin{itemize}
  \item perspektywy wielopodmiotowej obejmującej cały łańcuch ICT, 
  \item oceny opartej na weryfikowalnych dowodach jako mechanizmu przeciwdziałającego deklaratywnej zgodności,
  \item wieloperspektywicznej analizy uwzględniającej punkt widzenia prawnika, menedżera biznesowego i~inżyniera IT, \item rozszerzalności sektorowej pozwalającej na adaptację do specyfiki regulacyjnej różnych branż przy zachowaniu wspólnego rdzenia metodycznego.
\end{itemize}

Rezultatem jest rozdrobniona praktyka, w której oceny prawne, techniczne i~biznesowe są przeprowadzane w izolacji, podwykonawcy pozostają niesprawdzani, a wnioski z audytu opierają się na deklaracjach, a nie na weryfikowalnych dowodach.
Audyty przeprowadzane według istniejących metodyk często kończą się certyfikatami potwierdzającymi zgodność procesową, które nie odzwierciedlają faktycznego stanu bezpieczeństwa wdrożeń. Zjawisko to, zwane  \textit{compliance theater}, wynika ze słabości systemów oceny opartych na deklaracjach.  Organizacja może, z pełną znajomością swoich procedur oświadczyć, że wdrożyła wymagane środki ochrony, podczas gdy stan faktyczny jest odmienny, dokumentacja jest nieaktualna, mechanizmy kontrolne nieskuteczne, a dostawcy zewnętrzni nie podlegają żadnej weryfikacji.

\section{Postulaty i~cele ORION}

Metodologia ORION (\textit{Ocena Ryzyka Nadzorowanych Systemów Informatycznych}) nie ma na celu zastępowania metodologii audytowania zgodnych z ISO czy NIST. Wypełnia istotną lukę
w analizie wdrożeń systemów informatycznych z punktu widzenia zgodności z regulacjami. Jest to system oparty na dowodach, ponieważ, jak wykazują badania przytoczone w rozdziale
poświęconym walidacji odpowiedzi, organizacje nagminnie deklarują zgodność, której nie są w stanie udokumentować, a ta funkcja jest kluczowa.

Audyt prowadzony w ORION jest wykonywany ogólnie dla przedsiębiorstwa jako całości, analizując jego zgodność prawną z wymogami regulacyjnymi, a następnie,  zgodnie z przejściem grafu ICT, przeprowadzany jest na wszystkich systemach i~ich dostawcach. Identyfikowani są wszyscy
poddostawcy w łańcuchu ICT i~wszyscy podlegają badaniu. Wynikiem działania systemu jest określenie, który dostawca jest kluczowy, oraz wyznaczenie poziomu zgodności z prawem i~rekomendacji naprawczych z przypisanymi priorytetami.

ORION wprowadza cztery kluczowe innowacje metodyczne względem istniejących frameworków.

Pierwszą z nich jest zasada wieloperspektywiczności. ORION przyjmuje, że każdy element składowy infrastruktury ICT (system informatyczny, dostawca, wymóg, etc), jest rozumiany odmiennie przez prawnika, menedżera biznesowego, architekta oprogramowania czy administratora infrastruktury. Odmienność ta nie jest wadą, którą należy eliminować, jest naturalną cechą wynikającą z różnych wyzwań i~zadań na poszczególnych stanowiskach, którą metodyka musi uwzględniać. Perspektywa nie jest wyłącznie kategorią opisową, determinuje, kto w organizacji jest właściwą osobą do udzielenia odpowiedzi na konkretne pytanie audytowe, co eliminuje jeden z fundamentalnych błędów istniejącej praktyki: odpowiadanie na pytania prawne przez inżynierów i~na pytania techniczne przez prawników.

Drugą innowacją jest ocena oparta na dowodach. Każda odpowiedź na pytanie audytowe podlega weryfikacji przez pryzmat dostarczonego dowodu materialnego. Metodyka rozróżnia trzy poziomy jakości dowodu: brak dowodu, deklarację pisemną oraz dowód potwierdzony materialnie: log systemowy, certyfikat, wynik testu penetracyjnego, protokół z próby odtworzenia systemu.  Dodatkowo, dowód pisemny jest traktowany w zależności od jego sporządzenia, inaczej system bierze pod uwagę oświadczenie pisemne czytelnie podpisane, inaczej oświadczenie pisemne będące wynikiem anonimowego dokumentu. Brak dowodu nie zmienia odpowiedzi twierdzącej na przeczącą, lecz obniża jej efektywną wartość w modelu scoringowym w sposób kontrolowany i~konfigurowalny. Mechanizm ten precyzyjnie odzwierciedla rzeczywistość audytu: organizacja może  stosować właściwe praktyki bez ich pełnej dokumentacji, ale twierdzenie bez weryfikacji ma mniejszą wartość poznawczą niż twierdzenie udokumentowane. Jednoczenie konieczność przygotowania dowodu, co więcej, osobistego podpisania się pod dowodem, sktutjuje wysoką rzetelnością odpowiedzi. Jest to bezpośrednia odpowiedź na compliance theater.

Trzecią innowacją jest grafowe modelowanie łańcucha ICT z propagacją ryzyka. Audyt w ORION nie jest przeprowadzany wyłącznie dla podmiotu nadzorowanego traktowanego jako niepodzielna całość, obejmuje wszystkich dostawców i~podddostawców ICT, reprezentowanych jako wierzchołki skierowanego grafu zależności. Mechanizm systematycznego przejścia grafu zapewnia pełną identyfikację podmiotów w łańcuchu dostaw, które w tradycyjnych audytach pozostają poza polem widzenia. Algorytm propagacji ryzyka umożliwia z kolei obliczenie skumulowanego wpływu niezgodności na poziomie podddostawcy na ryzyko całego podmiotu nadzorowanego ,co odpowiada na pytanie, które organizacje zadają coraz częściej: czy mój dostawca jest bezpieczny, a jeśli nie, jak bardzo to mnie dotyczy?

Czwartą innowacją jest modułowa rozszerzalność sektorowa. System ORION jest zaprojektowany jako uniwersalne narzędzie z możliwością rozszerzenia na kolejne sektory. Rdzeń metodyczny, ontologia pojęć, model grafowy, formuły scoringu, struktura pytań i~raportowania, jest niezależny od branży i~może być stosowany wobec dowolnego podmiotu przetwarzającego dane w systemach ICT. Rozszerzenia domenowe, aktywowane w zależności od sektora audytowanego podmiotu, wprowadzają dodatkowe perspektywy, grupy pytań i~wymogi regulacyjne specyficzne dla danej branży. Oznacza to w praktyce, że ten system ORION może obsłużyć audyt banku, szpitala i~operatora energetycznego, różniąc się zestawem aktywnych pytań i~perspektyw, lecz zachowując identyczny model matematyczny, ontologię i~strukturę raportu. Ta cecha ORION ma istotne znaczenie nie tylko praktyczne, lecz także naukowe: pozwala na prowadzenie badań porównawczych między sektorami z użyciem jednolitej metodyki.



\section{Identyfikacja podmiotu audytowanego}
\label{sec:kim-jestem}

Przed przystąpieniem do właściwego audytu systemów informatycznych
metodologia ORION wymaga przeprowadzenia wstępnej klasyfikacji
podmiotu audytowanego. Celem tej fazy jest ustalenie, jakim
reżimem regulacyjnym podlega organizacja, w szczególności,
czy jest ona podmiotem kluczowym lub ważnym w rozumieniu dyrektywy NIS2~\autocite{dyrektywa_nis2}, podmiotem objętym zakresem stosowania rozporządzenia DORA~\autocite{dora2022}, czy też
jednostką podlegającą wyłącznie wymogom ogólnym, takim jak RODO~\autocite{rodo}. Wynik tej klasyfikacji determinuje, które moduły pytań zostaną aktywowane w dalszych etapach
audytu, jakie mnożniki wagowe zostaną przypisane do poszczególnych wymogów oraz które elementy, takie jak
plan ciągłości działania czy exit plan, mają charakter obowiązkowy, a które stanowią dobrowolne zalecenie.


Klasyfikacja podmiotu jest przeprowadzana metodą walidacji krzyżowej, gdzie zagadnienie analizowane jest niezależnie w perspektywie
zgodności, perspektywie użytkownika oraz perspektywie technicznej, a wynik jest ustalany na podstawie konsensusu odpowiedzi.
Wynikiem klasyfikacji jest przypisanie podmiotu do jednej z czterech kategorii regulacyjnych: \texttt{FINANSOWY} (podmioty objęte rozporządzeniem DORA), \texttt{KLUCZOWY} (podmioty kluczowe w rozumieniu dyrektywy NIS2 i~ustawy KSC), \texttt{WAŻNY} (podmioty ważne w rozumieniu NIS2 i~KSC) lub \texttt{POZOSTAŁE} (podmioty podlegające wyłącznie wymogom ogólnym, w~szczególności RODO). Kategoria determinuje zakres obowiązkowych elementów audytu oraz wartość mnożnika regulacyjnego $w_{reg}$, których szczegółową charakterystykę przedstawia tabela~\ref{tab:klasyfikacja-regulacyjna}.
% ============================================================
\section{Mechanizm walidacji jakości dowodów i~odpowiedzi}

\subsection{Struktura pytań}

Pytania stawiane w audycie mają z definicji wartość binarną. Oznacza to, że ankietowany moze odpowiedzieć na pytanie tylko i~wyłącznie \texttt{TAK} lub \texttt{NIE}. Każdemu pytaniu przypisana jest \texttt{ODPOWIEDŹ OCZEKIWANA},
rozumiana jako wartość wskazująca na stan zgodności z wymaganiem zgodne z semantyką treści pytania.

Działanie wartości oczekiwanej idealnie zilustruje przykład dwóch pytań:
\begin{center}
  Czy dwa plus dwa jest cztery?
\end{center}
Oraz pytanie drugie:
\begin{center}
  Czy dwa plus dwa różne od czterech?
\end{center}
W przypadku pierwszym odpowiedzią oczekiwaną jest \texttt{TAK}, w przypadku drugim  \texttt{NIE}. Współczynnik wagowy przypisywany jest wyłącznie wtedy, gdy odpowiedź jest zgodna z \texttt{ODPOWIEDZIĄ OCZEKIWANĄ}.

Zrezygnowano z opcji \texttt{NIE WIEM / BRAK DANYCH}. W sytuacji braku znajomości odpowiedzi oczekuje się, że ankietowany podejmie wysiłek uzyskania informacji. Brak pewności w odpowiedzi powinien
być rozstrzygany na niekorzyść audytu i~wynika to z konieczności ochrony organizacji.

Decyzja o rezygnacji z opcji \texttt{NIE WIEM / BRAK DANYCH} jest uzasadniona zarówno metodologicznie, jak i~z perspektywy zarządzania ryzykiem. Badania nad konstrukcją kwestionariuszy wykazują, że obecność
opcji neutralnej skłania respondentów do udzielania odpowiedzi zastępczych zamiast podejmowania wysiłku poznawczego niezbędnego do udzielenia odpowiedzi merytorycznej, zjawisko to określane jest w literaturze jako \textit{satisficing}~\autocite{Krosnick2002DontKnow}.
Kwestionariusze z wymuszonym wyborem binarnym generują dane wyższej jakości i~lepiej różnicują rzeczywiste stany wiedzy respondentów niż skale z opcją neutralną~\autocite{Dolnicar2011ForcedBinary}. 
W kontekście audytu bezpieczeństwa informatycznego brak wiedzy o stanie systemu jest sam w sobie wskaźnikiem ryzyka. Organizacja, która nie jest w stanie potwierdzić zgodności z wymaganiem, nie spełnia go z punktu
widzenia ochrony~\autocite{Waters2021DontKnow}. Przyjęcie zasady rozstrzygania wątpliwości na niekorzyść audytu odpowiada powszechnie stosowanej w zarządzaniu ryzykiem zasadzie, zgodnie z którą niepewność nie redukuje ryzyka, lecz je potęguje.

Dodatkowo, wprowadzona została opcja \texttt{flagi blokującej} $\text{FLAG}_{block}(\mathcal{S}) = \texttt{TRUE}$
 do pytań, których niespełnienie (odpowiedź inna niż \texttt{OCZEKIWANA}) oznacza przypisanie wartości zerowej dla całej sekcji. Są to pytania kluczowe, które muszą być bezwględnie spełnione.


\subsection{Model dowodowy}
\label{sec:dowody}


Każda odpowiedź na pytanie audytowe w ORION musi być opatrzona dowodem. Dowód nie zmienia wartości
odpowiedzi z \texttt{TAK} na \texttt{NIE} ani odwrotnie, lecz modyfikuje wagę odpowiedzi w modelu scoringowym.
Logika ta odzwierciedla cel audytu,  organizacja może faktycznie stosować właściwe praktyki bez ich pełnej i~faktycznej dokumentacji istnieje ryzyko, że dowód jest wadliwiwy.

Równie istotny jest aspekt behawioralny: badania z zakresu psychologii decyzji wskazują, że osoby podpisujące oświadczenia własnym imieniem i~nazwiskiem
kłamią istotnie rzadziej niż autorzy anonimowych dokumentów~\autocite{friesen2013designing}. Wymóg osobistego
podpisania dowodu jest zatem nie tylko formalny, lecz pełni funkcję mechanizmu antydefraudacyjnego wbudowanego w strukturę audytu.


\subsection{Walidacja odpowiedzi}

Walidacja odpowiedzi jest kluczowa, do prawidłowego przeprowadzenia procesu audytowania. Wskazują na to wyniki badań, które wskazują, że awersja do kłamstwa pojawia się głównie przy obowiązkowym raportowaniu \autocite{friesen2013designing}.

Osoby wypełniające dane ograniczają kłamstwo w wypełnianych ankietach, gdy są zobligowane do złożenia oświadczenia imiennego, potwierdzającego zgodne z prawdą wypełnienie ankiety. Badania wskazują, że współczynnik nieprawdziwych odpowiedzi wynosi \textit{g = 0.24}\autocite{peer2023commitment}. Współczynnik \textit{g} (zwany \textit{Hedges' g}) to standaryzowana miara wielkości efektu,  wyraża jak duża jest różnica między grupami w jednostkach odchylenia standardowego.


\subsection{Skala jakości dowodów}

Metodologia ORION wymaga przeprowadzenia odpowiedniej walidacji dowodu i~jakość dowodu, co bezpośrednio wpływa na wagę odpowiedzi.
ORION wyróżnia cztery poziomy jakości dowodu,zdefiniowane w tabeli \ref{tab:dowody}.
oznaczane symbolem $d \in \{0;\ 0{,}5;\ 0{,}8;\ 1{,}0\}$.


\begin{table}[H]
\label{sec:skala-dowodow}
\centering
\caption{Skala jakości dowodów w metodologii ORION}
\label{tab:dowody}
\renewcommand{\arraystretch}{1.6}
\begin{tabular}{p{1.5cm}p{2.5cm}p{4.5cm}p{4.5cm}}
\toprule
\textbf{Poziom} & \textbf{Symbol $d$}
& \textbf{Definicja}
& \textbf{Przykłady} \\
\midrule
Brak dowodu & $d = 0$
& Odpowiedź udzielona wyłącznie ustnie lub pisemnie
  bez jakiegokolwiek dokumentu potwierdzającego;
  niemożliwa do weryfikacji przez audytora.
& Oświadczenie słowne podczas wywiadu,
  odpowiedź wpisana bezpośrednio w formularz
  bez załącznika \\
\midrule
Dowód słaby & $d = 0{,}4$
& Dokument pisemny bez wskazania autora lub
  anonimowy; dokument ogólny niepowiązany
  wprost z audytowanym systemem.
& Regulamin ogólny dostawcy, polityka
  bezpieczeństwa bez daty i~podpisu,
  screen z systemu bez kontekstu \\
\midrule
Dowód średni & $d = 0{,}7$
& Oświadczenie pisemne z czytelnym podpisem
  osoby fizycznej i~datą; dokument wskazujący
  jednoznacznie na audytowany system.
& Podpisana deklaracja zgodności,
  umowa powierzenia z klauzulą TIA,
  certyfikat ISO z numerem seryjnym \\
\midrule
Dowód silny & $d = 1{,}0$
& Dowód materialny niemożliwy do sfałszowania
  bez pozostawienia śladu technicznego;
  log systemowy, wynik testu automatycznego,
  lub oświadczenie podpisane z poświadczeniem
  tożsamości wskazujące na pytanie audytowe.
& Log z systemu SIEM, wynik testu
  penetracyjnego, protokół z próby odtworzenia
  z BCP, oświadczenie potwierdzone podpisem osobistym. \\
\bottomrule
\end{tabular}
\end{table}

\subsection{Wpływ dowodu na wartość odpowiedzi}
\label{sec:wplyw-dowodu}

Efektywna wartość odpowiedzi na pytanie audytowe
jest iloczynem odpowiedzi binarnej i~jakości dowodu:

\begin{equation}
V_{eff}(q_j, \pi_i) = a_{ij} \cdot d_{ij}
\label{eq:veff}
\end{equation}

\noindent gdzie $a_{ij} \in \{0, 1\}$ jest odpowiedzią
perspektywy $\pi_i$ na pytanie $q_j$, a
$d_{ij} \in \{0;\ 0{,}4;\ 0{,}7;\ 1{,}0\}$
jest jakością dostarczonego dowodu.

Konsekwencje praktyczne wzoru~\eqref{eq:veff}
są następujące:
\begin{itemize}
  \item odpowiedź TAK bez dowodu ($a=1$, $d=0$)
  daje $V_{eff} = 0$ jest traktowana identycznie jak odpowiedź NIE,
  \item odpowiedź TAK z dowodem silnym ($a=1$, $d=1{,}0$)
  daje $V_{eff} = 1{,}0$  pełna wartość,
  \item odpowiedź NIE z dowodem silnym ($a=0$, $d=1{,}0$) daje $V_{eff} = 0$ wskazuje na niezgodność między odpowiedzią, a dowodem.
\end{itemize}

Mechanizm ten precyzyjnie odzwierciedla cel audytu: nie interesuje nas, co organizacja mówi, że robi, interesuje nas, co jest w stanie udowodnić.



\subsection{Walidacja krzyżowa oraz definicja konsensusu}
\label{sec:crossvalidation}
Walidacja krzyżowa (\textit{cross-validation}) jest kluczowym mechanizmem metodologii ORION, zapewniającym rzetelność klasyfikacji atrybutów i~odpowiedzi na pytania audytowe. To samo zagadnienie przedstawione jest użytkownikom z co najmniej dwóch perspektyw z trzech, a wynik jest ustalany na podstawie porównania odpowiedzi.
Mechanizm ten realizuje dwa cele jednocześnie: po pierwsze, wykrywa rozbieżności między perspektywami, które są
same w sobie wartościowymi sygnałami audytowymi; po drugie, zwiększa odporność systemu na celowe lub niecelowe podawanie nieprawdziwych informacji przez jeden z aktorów.


\begin{definition}[Formalna definicja konsensusu]
\label{sec:konsensus}
Niech $Q$ będzie zagadnieniem audytowym zadanym
perspektywom ze zbioru $\Pi_Q \subseteq \Pi$,
gdzie $|\Pi_Q| \geq 2$. Każda perspektywa
$\pi_i \in \Pi_Q$ udziela odpowiedzi
$a_i \in \{0, 1\}$ (NIE / TAK). Wynik
cross-validation zagadnienia $Q$ definiuje się jako:

\begin{equation}
\text{CV}(Q) =
\begin{cases}
1 & \text{jeśli } \left|\{i : a_i = 1\}\right|
    > \left|\{i : a_i = 0\}\right| \\
0 & \text{jeśli } \left|\{i : a_i = 0\}\right|
    > \left|\{i : a_i = 1\}\right| \\
\varnothing & \text{jeśli remis lub brak odpowiedzi}
\end{cases}
\label{eq:cv}
\end{equation}

Gdy $\text{CV}(Q) = \varnothing$, zagadnienie jest
oznaczane flagą \textsc{Rozbieżność} i~nie może
zostać zatwierdzone bez interwencji audytora.
W takim przypadku efektywna wartość zagadnienia
jest penalizowana:

\begin{equation}
V_{eff}(Q) = \min_{i}(a_i) \cdot \beta_{cv}
\label{eq:veff-rozbiez}
\end{equation}

\noindent gdzie $\beta_{cv} \in [0,1]$ jest
konfigurowalnym współczynnikiem kary za brak
konsensusu (wartość domyślna: $\beta_{cv} = 0{,}5$).

\end{definition}
Innymi słowy: brak zgody między perspektywami
jest traktowany jako połowicznie negatywna odpowiedź,
nawet jeśli jeden z aktorów odpowiedział twierdząco.

Gdy konsensus jest osiągnięty i~wszystkie perspektywy
dostarczyły dowody najwyższej jakości, efektywna wartość
jest dodatkowo premiowana:

\begin{equation}
V_{eff}(Q) = \text{CV}(Q) \cdot
\Bigl(1 + \gamma \cdot
\mathbb{1}\bigl[\forall\, i~: d_i = d_{max}\bigr]\Bigr)
\label{eq:veff-premia}
\end{equation}

\noindent gdzie $\gamma = 0{,}1$ jest premią za
pełny konsensus z kompletem dowodów najwyższej jakości,
a $d_{max}$ jest najwyższym poziomem dowodu
w skali opisanej w sekcji~\ref{sec:dowody}.

\subsection{Interpretacja rozbieżności}
\label{sec:interpretacja-rozbiez}

Rozbieżność między perspektywami nie jest błędem
procesu audytowego, jest wartościową informacją
merytoryczną. Tabela~\ref{tab:rozbiez-interpretacja}
przedstawia najczęstsze wzorce rozbieżności
i ich standardową interpretację w ORION.

\begin{table}[H]
\centering
\caption{Wzorce rozbieżności cross-validation
         i~ich interpretacja audytowa}
\label{tab:rozbiez-interpretacja}
\renewcommand{\arraystretch}{1.5}
\begin{tabular}{p{3.5cm}p{3.5cm}p{6cm}}
\toprule
\textbf{Wzorzec rozbieżności}
& \textbf{Typowa przyczyna}
& \textbf{Interpretacja i~zalecane działanie} \\
\midrule
Zgodność: TAK\newline Użytkownik: NIE\newline Techniczna: NIE
& Wymóg zdefiniowany w przepisach, lecz nieimplementowany
& Compliance theater: organizacja deklaruje zgodność formalną
  bez rzeczywistego wdrożenia. Flaga P1. \\
\midrule
Zgodność: NIE\newline Użytkownik: TAK\newline Techniczna: TAK
& Praktyka wdrożona bez formalnej podstawy prawnej
& Brak świadomości regulacyjnej. Ryzyko braku dokumentacji
  wymaganej przez regulatora. Rekomendacja formalizacji. \\
\midrule
Zgodność: TAK\newline Użytkownik: TAK\newline Techniczna: NIE
& Wymaganie znane, polityka istnieje, wdrożenie niezrealizowane
& Luka techniczna. Wymóg zdefiniowany lecz nieimplementowany
  na poziomie systemów. Flaga P2. \\
\midrule
Zgodność: NIE\newline Użytkownik: NIE\newline Techniczna: TAK
& System techniczny bez świadomości biznesowej i~prawnej
& Shadow IT lub wdrożenie bez akceptacji zarządczej.
  Wymaga formalnej inwentaryzacji i~decyzji. \\
\bottomrule
\end{tabular}
\end{table}



\section{Definicje i~aktorzy}

\subsection{Perspektywa}
\label{perspektywa-pytań}
Perspektywa punktu widzenia systemu określa interesariuszy uczestniczących w ocenie oraz sposób ich rozumienia analizowanych zagadnień. Każde pytanie ankietowe jest przypisane do co najmniej dwóch perspektyw, o ile jest to możliwe, do wszystkich trzech. Umożliwia to prowadzenie analizy porównawczej odpowiedzi na to samo kryterium, udzielonych niezależnie przez różnych interesariuszy. Rozbieżności między perspektywami stanowią sygnał do dalszego wyjaśnienia i~są równie wartościowe analitycznie jak odpowiedzi zgodne.

Zdefiniowane są trzy perspektywy:
\begin{enumerate}[label=(\roman*)]
  \item \textit{Perspektywa zgodności}, reprezentuje punkt widzenia osób odpowiedzialnych za przestrzeganie regulacji, takich jak prawnik, inspektor ochrony danych osobowych (IOD/DPO) czy specjalista ds. compliance. Celem tej perspektywy jest ocena kryterium pod kątem zgodności z obowiązującymi przepisami prawa, regulacjami rynkowymi oraz wewnętrznymi politykami organizacji (np. RODO, NIS2, DORA).

  \item  \textit{Perspektywa użytkownika}, reprezentuje jednostki organizacyjne firmy aktywnie korzystające z systemu w codziennej pracy, takie jak właściciel biznesowy systemu, kierownik działu lub kluczowy użytkownik końcowy. Celem tej perspektywy jest uzyskanie odpowiedzi odzwierciedlającej rzeczywisty sposób wykorzystania systemu, nie jego założenia projektowe, lecz faktyczny stan i~pola eksploatacji.

  \item \textit{Perspektywa techniczna}, reprezentuje punkt widzenia osób odpowiedzialnych za budowę, utrzymanie i~bezpieczeństwo systemu, takich jak architekt systemów, administrator lub inżynier. Celem tej perspektywy jest ocena, czy implementacja i~wdrożenie systemu zostały przeprowadzone zgodnie z obowiązującymi standardami technicznymi i~najlepszymi praktykami branżowymi.
\end{enumerate}


Przyjęte perspektywy determinują strukturę pytań ankietowych, w których pojedyncze kryterium analizowane jest przez pryzmat trzech różnych ujęć oraz wpływają na sposób definiowania obiektów. Ilustrując to podejście przykładem, zadajmy sobie pytanie:
\begin{center}
\vspace{1em} % Dodaje odstęp pionowy 
\textit{Jaka jest definicja systemu CRM?}
\end{center}


Na tak postawione pytanie zostaną udzielone trzy komplementarne odpowiedzi:

\begin{enumerate}[label=(\roman*)]
    \item Z punktu widzenia \textit{użytkownika}, system CRM służy do rejestrowania wszyskich kontaktów z klientami oraz realizacji procesów sprzedażowych.
    \item Z punktu widzenia \textit{zgodności}, system oraz jego fizyczna lokalizacja muszą gwarantować zgodność z polityką RODO, zapewniając organizacji pełną wiedzę o miejscu przetwarzania danych osobowych.
    \item Z punktu widzenia \textit{technicznego}, jest to rozwiązanie informatyczne, które ma zapewniać dostępność zasobów dla uprawnionych użytkowników poprzez określone kanały komunikacyjne.
\end{enumerate}


\subsection{Obiekt}

Centralnym pojęciem metodyki ORION jest \texttt{obiekt}, reprezentujący  byt w systemie ICT organizacji, który może być przedmiotem  oceny zgodności regulacyjnej. Pojęcie obiektu jest szerokie: obejmuje zarówno systemy informatyczne działające 
na infrastrukturze własnej organizacji, jak i~zewnętrznych dostawców usług, platformy chmurowe oraz poddostawców odkrytych w toku rozszerzania grafu ICT.

\begin{definition}[Obiekt]
Obiekt jest formalną jednostką audytową systemu ORION, 
reprezentującą byt w ekosystemie ICT organizacji, który:
\begin{enumerate}[label=(\roman*)]
    \item posiada tożsamość możliwą do jednoznacznej identyfikacji w grafie zależności ICT,
    \item może być przypisany do co najmniej jednego właściciela  lub podmiotu odpowiedzialnego,
    \item podlega co najmniej jednemu wymogowi regulacyjnemu wynikającemu z obowiązujących przepisów prawa,
    \item może być oceniany z co najmniej dwóch perspektyw audytowych zdefiniowanych w zbiorze $\Pi$.
\end{enumerate}
Formalnie, obiekt jest wierzchołkiem $v \in V$ grafu ICT 
$G = (V, E, W)$, posiadającym zbiór atrybutów  klasyfikacyjnych $\mathcal{A}(v)$ determinujących zakres i~intensywność przeprowadzanego audytu.
\end{definition}

W ORION obiektem jest zarówno aplikacja webowa obsługująca klientów banku, dostawca chmury publicznej, na której ta aplikacja działa, jak i~usługi firmy
 świadczącej usługi zarządzania kluczami kryptograficznymi na rzecz tego dostawcy chmury. Każdy z tych 
bytów jest odrębnym wierzchołkiem grafu, odrębnym podmiotem  audytu i~odrębnym źródłem ryzyka, choć wszystkie wymione podmioty uczestniczą w tym samym łańcuchu przetwarzania danych.

\subsubsection{Typy obiektów}

ORION wyróżnia cztery typy obiektów, które determinują domyślny zestaw aktywowanych sekcji pytań oraz rolę wierzchołka w grafie ICT.
Każdy obiekt posiada właściwości, czyli pola przyjmujące definiowane przez audytora wartości:

\begin{table}[H]
\centering
\caption{Typy obiektów w systemie ORION}
\label{tab:typy-artefaktow}
\renewcommand{\arraystretch}{1.5}
\begin{tabular}{p{3cm}p{7cm}p{5cm}}
\toprule
\textbf{Typ} & \textbf{Definicja} & \textbf{Przykłady} \\
\midrule
\texttt{PODMIOT} 
\newline Właściwości: nazwa
& Korzeń grafu ICT ($v_0$); audytowana organizacja jako całość. 
  Nie jest bezpośrednio oceniany pytaniami szczegółowymi, 
  jego wynik wynika z agregacji wszystkich pozostałych obiektów.
& Bank, szpital, operator infrastruktury krytycznej \\

\midrule
\texttt{SYSTEM\_INFORMATYCZNY}
\newline Właściwości: nazwa, dostawca, właściciel biznesowy
&  Oprogramowanie służące do gromadzenia, przetwarzania i~przesyłania informacji.
& System kadrowo-płacowy, baza danych pacjentów, 
  portal intranetowy, API płatności \\
\midrule

\texttt{DOSTAWCA} 
\newline Właściwości: nazwa, kraj pochodzenia, poddostawca
& Zewnętrzny podmiot prawny świadczący usługi ICT na rzecz 
  organizacji na podstawie umowy; odkryty na pierwszym 
  poziomie grafu.
& Dostawca ERP, operator chmury publicznej, 
  firma hostingowa, integrator systemów \\
\midrule
\texttt{PODDOSTAWCA}
\newline Właściwości: nazwa, kraj pochodzenia
& Zewnętrzny podmiot prawny świadczący usługi ICT 
  na rzecz dostawcy lub innego poddostawcy; 
  odkryty przez mechanizm triggera TIA lub deklarację 
  dostawcy wyższego poziomu.
& Google LLC jako operator infrastruktury dostawcy ERP, 
  Cloudflare jako CDN aplikacji webowej, 
  firma zarządzająca kluczami kryptograficznymi \\
\bottomrule
\end{tabular}
\end{table}

\subsection{Atrybuty obiektów}


Każdy obiektów jest scharakteryzowany przez zbiór atrybutów 
klasyfikacyjnych $\mathcal{A}(v)$, które pełnią dwie funkcje. 
Po pierwsze, determinują, które sekcje pytań audytowych 
zostaną aktywowane dla danego obiektów, na przykład 
system wewnętrzny działający wyłącznie on-premise nie 
otrzyma pytań dotyczących transferu 
danych do chmury. Po drugie, wyznaczają wagi stosowane 
w modelu scoringowym, decydując o tym, jak mocno wynik 
danego obiektu wpływa na wynik całego podmiotu.

\subsection{Atrybut: dostępność systemu}

Atrybuty obiektów są ustalane w fazie wstępnej audytu za pomocą pomocniczej ankiety, 
przed wygenerowaniem kwestionariusza szczegółowego. 
Widoczność systemu i~jego krytyczność, jest wyznaczana przez mechanizm  cross-validation z udziałem co najmniej dwóch perspektyw.


\begin{definition}[Poziom dostępności systemu]
Poziom dostępności obiektu $v$ jest wartością ze zbioru:
\[
\textit{dostępność}(v) \in 
\{\texttt{PRYWATNY},\ \texttt{PUBLICZNY}\}
\]
wyznaczaną przez mechanizm cross-validation na podstawie 
odpowiedzi co najmniej dwóch z trzech perspektyw rdzeniowych: 
\texttt{ZGODNOŚCI}, \texttt{UŻYTKOWNIKA} i~\texttt{TECHNICZNEJ}. 
\end{definition}


\begin{table}[H]
\centering
\caption{Poziomy dostępności obiektu, kryteria klasyfikacji
         i~mnożniki wagowe}
\label{tab:dostepnosc}
\renewcommand{\arraystretch}{1.5}
\begin{tabular}{p{2.8cm}p{1.5cm}p{4.5cm}p{4cm}}
\hline
\textbf{Poziom} & \textbf{$w_{krit}$} 
& \textbf{Kryteria klasyfikacji} 
& \textbf{Przykład} \\
\toprule
\texttt{PRYWATNY} & $1{,}0$ 
& System informatyczny działający na infrastrukturze 
  własnej organizacji lub zarządzanej wyłącznie przez nią; 
  bez bezpośredniego dostępu z sieci publicznej.
& System kadrowo-płacowy on-premise, baza danych pacjentów, 
  wewnętrzny portal intranetowy \\
\midrule
\texttt{PUBLICZNY} & $2{,}5$ 
& System informatyczny dostępny z sieci publicznej lub 
  świadczący usługi na rzecz podmiotów zewnętrznych; 
  obejmuje systemy SaaS, API i~aplikacje webowe.
& Portal klienta banku, system e-rejestracji, API płatności \\

\bottomrule
\end{tabular}
\end{table}

Klasyfikacja jest przeprowadzana przez cztery kryteria, z których każde jest zadawane wszystkim trzem perspektywom.
Logika wynikowa jest progowa: nawet jedna odpowiedź twierdząca z dowolnej perspektywy na dowolne kryterium
skutkuje klasyfikacją \texttt{PUBLICZNY}. Uzasadnienie tej asymetrii jest regulacyjne: jedna niepotwierdzona
ścieżka dostępu zewnętrznego jest wystarczającym powodem do traktowania systemu jako publicznie dostępnego
i stosowania względem niego pełnego zestawu wymogów bezpieczeństwa dla systemów publicznych.

\subsection{Atrybut: poziom krytyczności}

Poziom krytyczności jest atrybutem klasyfikacyjnym obiektu 
określającym, jak poważne skutki operacyjne, prawne i~społeczne 
pociągnęłaby za sobą niedostępność, naruszenie integralności 
lub utrata poufności danego systemu lub dostawcy. Atrybut ten 
jest wyznaczany w fazie wstępnej audytu metodą cross-validation, to samo zagadnienie jest zadawane niezależnie trzem perspektywom audytowym, a wynik klasyfikacji jest ustalany na podstawie 
konsensusu odpowiedzi.

\begin{definition}[Poziom krytyczności obiektu]
Poziom krytyczności obiektu $v$ jest wartością ze zbioru:
\[
\textit{krytyczność}(v) \in 
\{\texttt{KRYTYCZNY},\ \texttt{WAŻNY},\ \texttt{STANDARDOWY}\}
\]
wyznaczaną przez mechanizm cross-validation na podstawie 
odpowiedzi co najmniej dwóch z trzech perspektyw rdzeniowych: 
\texttt{ZGODNOŚCI}, \texttt{UŻYTKOWNIKA} i~\texttt{TECHNICZNEJ}. 
Poziom krytyczności determinuje mnożnik wagowy $w_{krit}(v)$ 
stosowany w równaniu agregacji wyników na poziomie 
podmiotu~\eqref{eq:Spodmiot_rozszerzone}.
\end{definition}

\begin{table}[H]
\centering
\caption{Poziomy krytyczności obiektu, kryteria klasyfikacji 
         i~mnożniki wagowe}
\label{tab:krytycznosc}
\renewcommand{\arraystretch}{1.5}
\begin{tabular}{p{2.8cm}p{1.5cm}p{4.5cm}p{4cm}}
\hline
\textbf{Poziom} & \textbf{$w_{krit}$} 
& \textbf{Kryteria klasyfikacji} 
& \textbf{Podstawa regulacyjna} \\
\toprule
\texttt{KRYTYCZNY} & $2{,}0$ 
& Niedostępność przez ponad 4~h powoduje bezpośrednie straty 
  operacyjne, naruszenie przepisów prawa lub zagrożenie 
  bezpieczeństwa osób; brak możliwości zastąpienia; 
  system wspiera funkcję kluczową w rozumieniu NIS2 lub DORA
& NIS2 Art.~3; DORA Art.~2; 
  ustawa o krajowym systemie cyberbezpieczeństwa \\
\midrule
\texttt{WAŻNY} & $1{,}5$ 
& Niedostępność poważnie utrudnia działalność, lecz jej 
  nie wstrzymuje; istnieje procedura zastępcza; 
  system wspiera funkcję ważną w rozumieniu NIS2
& NIS2 Art.~3 ust.~2; 
  ISO/IEC~22301 \\
\midrule
\texttt{STANDARDOWY} & $1{,}0$ 
& Niedostępność powoduje utrudnienia pomocnicze; 
  istnieje możliwość pracy bez systemu przez 
  co najmniej 48~h bez istotnych strat
& --- \\
\bottomrule
\end{tabular}
\end{table}

Klasyfikacja krytyczności jest przeprowadzana przez sześć kryteriów (K1--K6) zadawanych wszystkim trzem perspektywom. Wynik końcowy wyznaczany jest progowo: suma wszystkich odpowiedzi twierdzących ze wszystkich perspektyw i~wszystkich kryteriów jest normalizowana do skali klasyfikacyjnej zgodnie z~ankietą przedstawioną w~sekcji~\ref{sec:krok23}.

\subsection{Atrybut: klasyfikacja utrzymania w chmurze}
\label{sec:klasyfikacja-chmura}
Lokalizacja miejsca utrzymania usługi ma kluczowy wpływ na warunki jej świadczenia. Oprogramowanie działające \textit{on-premise}, czyli w siedzibie organizacji, nie jest narażone na ryzyko przesyłu danych do zewnętrznych dostawców. Istotna jest również prawna jurysdykcja podmiotu, który hostuje dane. Fakt przechowywania danych w chmurze będzie implikować konieczność dokładnego audytu łańcucha dostawców i~wykonania oceny skutków transferu (\textit{TIA}), wymaganej przez RODO w sytuacji, gdy dane przesyłane są z Europejskiego Obszaru Gospodarczego (\textit{EOG}) do tzw. "państwa trzeciego" (np. USA, Indie, Chiny).

\subsubsection*{Klasyfikacja utrzymania  w chmurze}
Metodologia ORION definiuje dwa typy lokalizacji danych, a ich atrybut wpływa na aktywacje sekcji pytań dotyczących rozwiązań chmurowych oraz transferu danych.


\begin{table}[H]
\centering
\caption{Klasyfikacja utrzymania w chmurze w ORION}
\label{tab:typy-chmura}
\renewcommand{\arraystretch}{1.4}
\begin{tabular}{p{5cm}p{8cm}}
\toprule
\textbf{Typ utrzymania} & \textbf{Opis}\\
\midrule
\texttt{CHMURA\_K3} & Dane przechowywane są poza organizacją u zewnętrznego podmiotu (SaaS, IaaS/PaaS, model hybrydowy, hosting, kolokacja), którego lokalizacja lub jurysdykcja nie gwarantuje przechowywania danych w EOG. Obejmuje dostawców podlegających prawu państw trzecich, np. Amazon Web Services, Google Cloud, Microsoft Azure.\\
\midrule
\texttt{CHMURA\_EOG} & Dane przechowywane są poza organizacją u zewnętrznego podmiotu (SaaS, IaaS/PaaS, model hybrydowy, hosting, kolokacja), którego lokalizacja i~jurysdykcja gwarantuje przechowywanie danych wyłącznie w EOG, np. NASK SA, Exea Data Center, Beyond.pl, OVH (region PL/EU).\\
\midrule
\texttt{ON\_PREMISE} & Dane przechowywane są wyłącznie na serwerach i~infrastrukturze będącej własnością organizacji, fizycznie zlokalizowanej w jej siedzibie lub własnym centrum danych.\\
\bottomrule
\end{tabular}
\end{table}

\subsection{Atrybut: klasyfikacja atrybutu typów przetwarzanych danych}
\label{sec:klasyfikacja-danych}

Dotychczasowe modele oceny ryzyka ICT nie uwzględniają charakter  przetwarzanych danych jako samodzielny wymiar ryzyka z formalnym  przełożeniem na wynik scoringowy. W metodologii ORION charakter danych jest traktowany jako atrybut pierwszorzędny, który wpływa bezpośrednio  na wagę artefaktu w agregacji końcowej. Logika jest prosta: naruszenie bezpieczeństwa systemu przetwarzającego tajemnicę lekarską lub dane finansowe klientów pociąga za sobą nieporównywalnie poważniejsze  skutki prawne, finansowe i~społeczne niż naruszenie systemu  obsługującego wyłącznie dane wewnętrzne organizacji.

\subsubsection{Klasyfikacja typów danych}

Metodologia ORION wyróżnia cztery typy przetwarzanych danych, uporządkowane  rosnąco według wrażliwości prawnej i~społecznej. Każdemu typowi  przypisany jest mnożnik wagowy $w_{data}$, który modyfikuje 
wpływ artefaktu na wynik podmiotu zgodnie 
z równaniem~\eqref{eq:Spodmiot_rozszerzone}.

\begin{table}[H]
\centering
\caption{Klasyfikacja typów danych w ORION i~odpowiadające mnożniki wagowe}
\label{tab:typy-danych}
\renewcommand{\arraystretch}{1.4}
\begin{tabular}{p{2.5cm}p{1.5cm}p{1.5cm}p{3.5cm}p{4.5cm}}
\toprule
\textbf{Typ danych} & \textbf{Symbol} & \textbf{$w_{data}$} 
& \textbf{Podstawa prawna} & \textbf{Przykłady} \\
\midrule
\texttt{BRAK} & $T_0$ & $1{,}0$ 
& --- 
& Dane konfiguracyjne, logi techniczne bez PII \\
\midrule
\texttt{OSOBOWE} & $T_1$ & $1{,}3$ 
& RODO Art.~4 ust.~1 
& Imię, adres, adres e-mail, numer IP \\
\midrule
\texttt{WRAŻLIWE} & $T_2$ & $1{,}6$ 
& RODO Art.~9 ust.~1 
& Stan zdrowia, dane biometryczne, 
  przekonania religijne, orientacja seksualna \\
\midrule
\texttt{TAJEMNICA} & $T_3$ & $2{,}0$ 
& Ustawa sektorowa\newline(bankowa, lekarska, adwokacka) 
& Tajemnica bankowa, tajemnica lekarska, 
  tajemnica zawodowa, dane niejawne \\
\bottomrule
\end{tabular}
\end{table}

\subsection{Atrybut: kraj pochodzenia}
\label{sec:kraj-pochodzenia}
Atrybut ten określa kraj pochodzenia podmiotu w celu jednoznacznego ustalenia jego statusu prawnego, tj. czy jest to podmiot mający siedzibę na terenie Europejskiego Obszaru Gospodarczego (EOG), czy też kwalifikuje się on jako podmiot z kraju trzeciego w myśl obowiązujących przepisów prawa (w szczególności w kontekście regulacji dotyczących ochrony danych osobowych RODO oraz cyberbezpieczeństwa).

Wartość tego atrybutu ma kluczowe znaczenie dla logiki warunkowego uruchamiania sekcji pytań podczas procesu audytowego. W przypadku zidentyfikowania podmiotu pochodzącego z kraju spoza EOG, system automatycznie rozszerza zakres audytu o dodatkowe wymagania. Wynika to z podwyższonego ryzyka operacyjnego i~prawnego (np. w obszarze transferu danych lub łańcucha dostaw). W konsekwencji na organizację audytującą nakładany jest obowiązek przeprowadzenia pogłębionej weryfikacji, obejmującej m.in. szczegółowe badanie mechanizmów zabezpieczeń, umów transferowych oraz analizę wpływu na ryzyko (np. TIA).


\begin{table}[H]
\centering % Opcjonalnie: wyśrodkowanie tabeli na stronie
\caption{Dopuszczalne wartości pochodzenia przedsiębiorstw }
\label{tab:country-values}

\resizebox{\textwidth}{!}{%
\begin{tabular}{p{5cm} p{11cm}}
\toprule
\textbf{Kod kraju} & \textbf{Opis i~zakres odpowiedzialności} \\
\midrule
\texttt{KOD ISO} & Kod kraju pochodzenia, np. PL, DE, US \\
\bottomrule
\end{tabular}%
}
\end{table}

\section{Plan ciągłości dzialania}

Plan ciągłości działania jest procedurą wymaganą od podmiotów kluczowych zarówno na gruncie ustawy o KSC\autocite{uksc2026}, jak i~rozporządzenia DORA\autocite{dora2022}.

\begin{definition}[Plan ciągłości działania (BCP)]
Zbiór udokumentowanych procedur, które określają działania, jakie podmiot musi podjąć w celu odzyskania zdolności do wykonywania operacji oraz zapewnienia ciągłości procesów krytycznych na akceptowalnym poziomie po wystąpieniu incydentu zakłócającego.
\end{definition}

Podstawą planu ciągłości działania jest określenie dwóch wskaźników odporności. Pierwszy, RTO (\textit{Recovery Time Objective}) określa maksymalny dopuszczalny czas, przez jaki system może pozostawać niedostępny po wystąpieniu awarii, im niższy, tym wyższe wymagania wobec infrastruktury. Drugi, RPO (Recovery\textit{ Point Objective}) wyznacza  maksymalną akceptowalną utratę danych, wyrażoną jako punkt w czasie, do którego dane muszą zostać odtworzone, ilustrując przykładem, RPO wynoszące 15 minut oznacza, że żadna transakcja starsza niż kwadrans nie może zostać bezpowrotnie utracona.

Plan musi przewidywać konkretne scenariusze postępowania na wypadek różnych klas awarii systemów IT, ataku ransomware czy niedostępności kluczowego personelu. Równolegle obowiązują procedury wykrywania, klasyfikacji i~eskalacji incydentów, w tym raportowania do organów nadzorczych, takich jak KNF, w ściśle określonych ramach czasowych wynikających z obiwązującego prawa, a ich definicje szczegółowo określa DORA i~KSC.

Niezbędne jest opracowanie polityki wytwarzania oraz przywracania kopii zapasowych (DRP). Kopie zapasowe muszą być fizycznie lub logicznie odizolowane od głównej infrastuktury produkcyjnej, co chroni jest przed zaszyfrowaniem w przypadku ataku ransomware. Plan odtwarzania precyzuje harmonogram tworzenia kopii, lokalizację ich przechowywania, odpowiedzialnych za proces oraz procedury testowego odtwarzania danych w kontrolowanych warunkach.

Bardzo ważnym elementem plan ciągłości działania jest strategia komunikacji na czas kryzysu, który określa, kto, kiedy i~w jaki sposób informuje klientów, pracowników, kontrahentów oraz nadzór o wystąpieniu awarii i~przewidywanym czasie jej usunięcia. 

Zgodnie z wymogami DORA, instytucja finansowa nie może ograniczyć planowania ciągłości wyłącznie do własnych zasobów, obowiązkiem jest zapewnienie, że dostawcy usług chmurowych i~IT również dysponują udokumentowanymi planami ciągłości i~są w stanie dotrzymać uzgodnionych poziomów usług (\textit{SLA}) w sytuacji kryzysowej. Wymóg ten powinien być wprost zawarty w umowach z dostawcami, wraz z prawem do audytu.

Plan ciągłości działania podlega obowiązkowym testom ich skuteczności, które obejmują symulacje cyberataków, przełączanie ruchu na ośrodek zapasowy oraz testy odtwarzania danych z kopii zapasowych. Wyniki każdego testu powinny skutkować aktualizacją dokumentacji.

Odpowiedzialność za zatwierdzenie i~skuteczność planów ciągłości działąnia, spoczywa na zarządzie instytucji. Od strony technicznej plan musi eliminować Single Point of Failure (\textit{SPOF}) poprzez zapewnienie redundantnych systemów, łączy komunikacyjnych i~zasilania, tak aby awaria pojedynczego komponentu nie prowadziła do niedostępności całej usługi.

W metodoligii ORION, plan ciągłości działania jest wymagany dla podmiotów określonych w KSC oraz DORA. Jego brak lub nieprawidłowe zdefiniowanie w organizacjach, gdzie jest on obowiązkowy wskazują binarnie na uchybienia w działaniu organizacji.

\begin{table}[H]
\centering
\caption{Mnożniki wagowe BCP wg typu organizacji}
\label{tab:bcp-wagi}
\renewcommand{\arraystretch}{1.4}
\begin{tabular}{p{4.5cm}p{1.5cm}p{2.5cm}p{4.5cm}}
\toprule
\textbf{Typ organizacji} & \textbf{Symbol} & \textbf{$w_{reg}$} & Obowiązkowe\\
\midrule
\texttt{FINANSOWY} & $D_0$ & $1{,}0$   & tak\\
\midrule
\texttt{KLUCZOWY} & $D_1$ & $1{,}0$   & tak\\
\midrule
\texttt{WAŻNY} & $D_2$ & $1{,}0$   & tak\\
\midrule
\texttt{POZOSTAŁE} & $D_3$ & $0{,}5$ & nie \\
\bottomrule
\end{tabular}
\end{table}

\section{Exit plan}

Exit plan (plan wyjścia) to obowiązkowy element zarządzania dostawcami IT w banku, wymagany przez rozporządzenie DORA oraz ustawę KSC.

\begin{definition}[Exit plan]
Strategia wyjścia z usługi pozwalająca organizacji wyłączyć system, zakończyć współpracę z dotychczasowym dostawcą oraz ponownie uruchomić usługę u innego dostawcy przy zachowaniu ciągłości działania lub minimalizacji przestoju.
\end{definition}

Plan wyjścia stanowi pisemnie udokumentowaną procedurę umożliwiającą organizacji zakończenie współpracy z dostawcą ICT i~przeniesienie usług do alternatywnego środowiska, bez przestoju operacyjnego i~bez utraty danych. Definicja ma zagwarantować ciągłości działania, czyli zapewnić minimalny czas niedostępności usługi dla jej użytkowników. Równie istotna jest ochrona przed zjawiskiem \textit{vendor lock-in}, czyli uzależnienia od jednego dostawcy

Wedługi definicji DORA, plan musi identyfikować musi wskazywać rzeczywistego dostawcę zastępczego lub opisywać scenariusz powrotu do infrastruktury własnej. Co istotne, umowa z dostawcą musi zawierać klauzulę zobowiązującą go do udostępnienia wszystkich danych w formacie umożliwiającym ich import do innego systemu

Dostawca ICT jest zobowiązany do świadczenia usług przez określony czas po wypowiedzeniu umowy, tak aby organizacja przeprowadziła migrację danych. Co istotne, wsparcie to nie może mieć charakteru pasywnego, dostawca ma obowiązek aktywnej pomocy technicznej przy przenoszeniu procesów, danych i~konfiguracji do nowego środowiska. Tego rodzaju zobowiązania muszą być precyzyjnie zapisane w umowie, z określonymi terminami i~zakresem odpowiedzialności.

Exit plan traci swoją wartość, jeśli nie jest regularnie weryfikowany. DORA wymaga przeprowadzania testów planu co najmniej raz w roku lub każdorazowo po istotnych zmianach w architekturze systemu (Art. 28 ust. 4\autocite{dora2022}), również ustawa KSC (Art. 8\autocite{uksc2026}) wymaga opracowania exit planu.  Całość procesu musi być zaplanowana tak, aby klient instytucji w ogóle nie odczuł zmiany dostawcy.

Ważma jest forma przygotowania dokumentacji. DORA, w art 28 ust. 4\autocite{dora2022} wymaga oddzielnego dokumentu, natomiast KSC wymaga realizuje cele exit planu poprzez klauzule umowne z dostawcą.

W metodoligii ORION, exit plan jest wymagany dla podmiotów określonych w KSC oraz DORA. Podobnie jak plan ciągłości działania,  brak lub nieprawidłowo zdefiniowany exit plan, gdu jest obowiązkowy, wskazuje binarnie na uchybienia w działaniu organizacji.

\begin{table}[H]
\centering
\caption{Mnożniki wagowe exit plan wg typu organizacji}
\label{tab:exitplan-wagi}
\renewcommand{\arraystretch}{1.4}
\begin{tabular}{p{4.5cm}p{1.5cm}p{2.5cm}p{4.5cm}}
\toprule
\textbf{Typ organizacji} & \textbf{Symbol} & \textbf{$w_{reg}$} & Obowiązkowe\\
\midrule
\texttt{FINANSOWY} & $D_0$ & $1{,}0$   & tak\\
\midrule
\texttt{KLUCZOWY} & $D_1$ & $1{,}0$   & tak\\
\midrule
\texttt{WAŻNY} & $D_2$ & $1{,}0$   & nie\\
\midrule
\texttt{POZOSTAŁE} & $D_3$ & $0{,}5$ & nie \\
\bottomrule

\end{tabular}
\end{table}




% \section{Analiza}

% Nąstepnie kazde pytanie zadane do kazdego systemu (artefaktu):

% Wymóg regulacyjny - czyli w czy wartość wynika z ustawy i~mysi być spełniona, czy też jest dobrą praktyką, więc opcjonalnie.
% Punkt widzenia artefaktu - prawne, biznesowe, techniczne w zakresie oprogramowania, techniczne w zakresie wdrożenia. 
% Dotyczy to samego pytania, ktore pomoze na wejsciu ocenic kto powien na pytanie odpowiedziec.

% i~odpowiedzi przypisywane są do binarnych wyników.

% Dowód - czy przeprowadzony zostal dowod. Nie jest on wymagany, ale jego brak obniza wartosc. Jest to wartosc wylisotwana: 0, 0.5 (jest deklaracja), 1 - potwierdzeni dowodem, ktora okresla jakosc dowodu. Waga dowodu - jesli inna niz 1 to rosnie ryzyka mocniej, zaproponuj jak.

% Waga ryzyka -  okreslnie specjalnej wagi ryzyka, domyslnie 1. Globalna waga dla pytania.

% Następnie kazda odpowiedz grupowana jest w grupy i~powstaje procentowy wynik na podstawie wktora wartosci uwzglednaijac wage. 



% Wnioskowanie:

% Ryzyko

% Poziom zgodności

% Luka zgodności

% Każde pojęcie:

% definicja słowna

% definicja formalna

% relacje z innymi

% \textbf{zrobie zgodnosc binarną}


